\documentclass[10pt,a4paper]{amsart}
\usepackage[margin=19mm]{geometry}
\usepackage{amsmath,amssymb,mathtools}
\usepackage[scr=rsfs,cal=euler]{mathalfa}
\usepackage{xfrac}
\usepackage[unicode,hidelinks,bookmarks=false]{hyperref}
\newcommand{\ie}{i.e.\ }
\newcommand{\cf}{cf.\ }
\newcommand{\resp}{respectively\ }
\newcommand{\Subsection}{\subsection}
\providecommand{\nobreakdash}{\nobreak\hbox{-}}
\setlength{\emergencystretch}{2em}
\multlinegap0pt
\usepackage{bm}
\usepackage{bbm}
\usepackage{dsfont}
\usepackage{xspace}
\usepackage{cancel}
\usepackage{mathtools}
\usepackage{amsthm}
\makeatletter
\@ifundefined{lemma}{\newtheorem{lemma}{Lemma}}{}
\makeatother
\title[On the Injectivity of STFT Phase Retrieval with super-expone]{On the Injectivity of STFT Phase Retrieval with super-exponential decaying window function}
\author{Shuang Guan, Kasso A. Okoudjou}
\date{Source: arXiv:2511.15894v2 (CC BY 4.0)}
\begin{document}
\maketitle

\noindent\emph{Source and licence.} On the Injectivity of STFT Phase Retrieval with super-exponential decaying window function. Version arXiv:2511.15894v2 (CC BY 4.0), distributed under the Creative Commons Attribution 4.0 International licence, as recorded in the arXiv licence metadata for this exact version and shown on the abstract page https://arxiv.org/abs/2511.15894v2. Full rights notice: licenses/arxiv-2511.15894-CC-BY-4.0.txt.\par\smallskip

\noindent\textit{Editorial scope.} Counted unit: the Lemma whose source label is \texttt{OrderType} in the article (source0.tex lines 210--212, source label \texttt{OrderType}), delivered together with the article's own complete original proof of it (source0.tex lines 213--280, one \texttt{proof} environment of 68 lines). No mathematics is written, repaired or re-derived: statement and proof are the article's own text line for line. Required context retained uncounted: TypeCalculation (lines 204--206). Every definition, display and label of those lines is retained unchanged. Omitted from this package and not redistributed: 1--203, 207--209, 281--717. Nothing inside a retained excerpt is omitted or altered: every retained body is delivered exactly as the article's lines are written, so no omission receipt of any kind accompanies this package. Citations: the retained excerpts carry no citation command at all, so no bibliography entry of the article is transcribed and no reference is redistributed. Reference adaptation: none was needed and none was made; every reference inside the delivered unit resolves inside it, so no unresolved reference or citation is produced. Printed slips kept literal: no printed word, symbol, hypothesis or number of the counted statement or of its proof was altered. Layout and class: the article's own class is \texttt{svmult} with packages including type1cm, makeidx, graphicx, multicol, footmisc, newtxtext, newtxmath; the wrapper is the lane's pinned \texttt{amsart} editorial wrapper at 10pt on A4, which restates the theorem-like environments the retained text needs and adds the article's own macro definitions that the retained text needs (none), with extra wrapper packages (bm, bbm, dsfont, xspace, cancel, mathtools). The delivered numbering is the unit's own. Assets: no retained span contains a figure environment, a graphics-inclusion command, a PDF-page inclusion, a table or a TikZ picture; no image file is copied into this package, the package declares no asset dependency and no image file is redistributed with it. Licence: version arXiv:2511.15894v2 of this article is distributed under the Creative Commons Attribution 4.0 International licence, as recorded in the arXiv licence metadata for this exact version and shown on the abstract page https://arxiv.org/abs/2511.15894v2. A full rights notice is delivered at licenses/arxiv-2511.15894-CC-BY-4.0.txt. Characters: every retained excerpt is pure ASCII, so the delivered engine, XeLaTeX with its default Latin Modern text font, reports no missing character.
    \begin{equation}\label{TypeCalculation}
        \tau = \frac{1}{e\rho} \limsup_{n \to\infty} (n \sqrt[n]{|c_n|^{\rho}}).
    \end{equation}

\begin{lemma}\label{OrderType}
Let $f \in L^2(\mathbb{R})$ be a function whose Fourier transform satisfies $|\hat{f}(\xi)| \le C e^{-a|\xi|^m}$ for some $m>1$ and $a>0$. Then its analytic continuation $f(z)$ is an entire function of type at most $\tau = \frac{m-1}{m}(2\pi)^{m/(m-1)}(am)^{-1/(m-1)}$ corresponding to maximum possible order $\rho = \frac{m}{m-1}$.
\end{lemma}

\begin{proof}
The analytic continuation of the function, $f(z)$, is given by the inverse Fourier transform:
\begin{equation*}
    f(z) = \int_{-\infty}^\infty \hat{f}(\xi) e^{2\pi i z \xi} d\xi
\end{equation*}
The Taylor coefficients of $f(z)$ are given by 
%$c_n = \frac{f^{(n)}(0)}{n!}$\textcolor{red}{, which can be found as}
\begin{equation*}
c_n = \frac{f^{(n)}(0)}{n!} = \frac{(2\pi i)^n}{n!} \int_{-\infty}^\infty \xi^n \hat{f}(\xi) d\xi
\end{equation*}
We can now bound the magnitude of the coefficients using the given decay of $\hat{f}(\xi)$:
\begin{equation*}
|c_n| \le \frac{(2\pi)^n}{n!} \int_{-\infty}^\infty |\xi|^n |\hat{f}(\xi)| d\xi \le \frac{C(2\pi)^n}{n!} \int_{-\infty}^\infty |\xi|^n e^{-a|\xi|^m} d\xi
\end{equation*}
The integral can be evaluated using the Gamma function:
\begin{equation*}
    \int_{-\infty}^\infty |\xi|^n e^{-a|\xi|^m} d\xi = \frac{2}{m a^{(n+1)/m}} \Gamma\left(\frac{n+1}{m}\right). 
\end{equation*}
Using Stirling's formula for Gamma function, we obtain 
\begin{equation*}
 \Gamma(\frac{n+1}{m}) = \sqrt{\frac{2 \pi m}{n+1}}(\frac{n+1}{me})^{(n+1)/m}\big(1+ O(\frac{m}{n+1})\big).
\end{equation*}
and 
\begin{equation*}
    n! \approx \sqrt{2 \pi n}(\frac{n}{e})^n
\end{equation*}
Then, for some $C' >0$, we derive the following expression:
         \begin{align}\label{GeneralTaylorCoefficient}
              |c_n|
             &\lesssim \Big((2\pi)^n \cdot n^{-\frac{1}{2}-n}\cdot e^{n -\frac{n+1}{m}}\cdot a^{-\frac{n+1}{m}}\cdot m ^{-\frac{1}{2}-\frac{n+1}{m}}\cdot (n+1)^{\frac{n+1}{m}-\frac{1}{2}}\Big)\cdot\Big(1 + \frac{C'm}{n+1}\Big) \\
               &=  R(a,n,m)\cdot n^{-n}\cdot(n+1)^{\frac{n+1}{m}}\cdot\Big(1 + \frac{C'm}{n+1}\Big) \nonumber
        \end{align}
The term $R(a,n,m)$ collects all remaining exponential and polynomial factors. It follows that:
\begin{align*}
    \rho &= \limsup_{n\to\infty} \frac{n \log n}{-\log|c_n|} \lesssim \limsup_{n\to\infty} \frac{n\log n}{-\log R(a,m,n) + n\log n -\frac{n+1}{m}\log(n+1)- \log(1+\frac{C'm}{n+1})}\\
    &\lesssim \limsup_{n\to\infty} \frac{1}{\frac{-\log R(a,m,n)}{n\log n} + 1 -\frac{n+1}{nm}\frac{\log(n+1)}{\log n}- \frac{\log(1+\frac{C'm}{n+1})}{n\log n}}.
\end{align*}
A series of computations shows that $$\begin{cases} \lim_{n\to \infty}\frac{\log R(a,m,n)}{n\log n} =0\\ \lim_{n\to \infty}\frac{\log(1+ \frac{C'm}{n+1})}{n\log n}=0\end{cases}.$$  Therefore, the remaining terms are:
\begin{equation*}   
    \rho \leq \limsup_{n\to \infty}\frac{1}{1-\frac{n+1}{mn}\frac{\log(n+1)}{\log n}}= \frac{1}{1- \frac{1}{m}} = \frac{m}{m-1}.
\end{equation*}
Let $\tilde{\rho} = \frac{m}{m-1}$, and choose $A,R >0$ such that 
\begin{equation}\label{TypeWRTorder}
    \max_{|z|= r}|f(z)| < Re^{Ar^{\tilde{\rho}}}.    
\end{equation}
Now we compute the type $\tau$ of $f(z)$ with respect to order $\tilde{\rho}$, i.e., the greatest lower bound of $A$ in Equation~\eqref{TypeWRTorder}. We note that 

$$ |c_n|^{\tilde{\rho}/n} \lesssim \Big(n^{-\frac{1}{2}-n}(n+1)^{\frac{n+1}{m}-\frac{1}{2}}\Big)^{\tilde{\rho}/n}\, (e^{n - \frac{n+1}{m}})^{\tilde{\rho}/n}\, a^{-\frac{n+1}{m}\cdot \frac{\tilde{\rho}}{n}} \, m ^{(-\frac{1}{2}-\frac{n+1}{m})\cdot \frac{\tilde{\rho}}{n}}\, \Big(1 + \frac{C'm}{n+1}\Big)^{\tilde{\rho}/n} $$ But,

$$\Big(n^{-\frac{1}{2}-n}(n+1)^{\frac{n+1}{m}-\frac{1}{2}}\Big)^{\tilde{\rho}/n} = n^{-\frac{\tilde{\rho}}{2n}-\tilde{\rho}}(n+1)^{\frac{(n+1)\tilde{\rho}}{mn}-\frac{\tilde{\rho}}{2n}}=\Big(n^{-\tilde{\rho}}(n+1)^{\frac{(n+1)\tilde{\rho}}{mn}}\Big)\Big( n(n+1) \Big)^{^{-\frac{\tilde{\rho}}{2n}}}$$
Thus,

$$n|c_n|^{\tilde{\rho}/n}\lesssim  \Big(n^{1-\tilde{\rho}}(n+1)^{\frac{(n+1)\tilde{\rho}}{mn}}\Big)\Big( n(n+1) \Big)^{^{-\frac{\tilde{\rho}}{2n}}}(e^{n - \frac{n+1}{m}})^{\tilde{\rho}/n}\, a^{-\frac{n+1}{m}\cdot \frac{\tilde{\rho}}{n}} \, m ^{(-\frac{1}{2}-\frac{n+1}{m})\cdot \frac{\tilde{\rho}}{n}}\Big(1 + \frac{C'm}{n+1}\Big)^{\tilde{\rho}/n}$$

However, we see that 
$$\begin{cases}  \lim_{n\to \infty}(e^{n - \frac{n+1}{m}})^{\tilde{\rho}/n} = e^{\tilde{\rho}- \frac{\tilde{\rho}}{m}} = e^{\frac{m}{m-1} - \frac{1}{m-1}}=e\\
\lim_{n\to \infty} a^{-\frac{n+1}{m}\cdot \frac{\tilde{\rho}}{n}} = a^{-\tilde{\rho}/m}\\\lim_{n\to\infty}m ^{(-\frac{1}{2}-\frac{n+1}{m})\cdot \frac{\tilde{\rho}}{n}} = m^{-\tilde{\rho}/m}\\
\lim_{n\to \infty}\Big(1 + \frac{C'm}{n+1}\Big)^{\tilde{\rho}/n} =1\end{cases}.$$ Moreover,


$$\lim_{n\to \infty}\Big( n(n+1) \Big)^{^{-\frac{\tilde{\rho}}{2n}}} = 1 \;\quad \text{and} \lim_{n\to \infty}\Big(n^{1-\tilde{\rho}}(n+1)^{\frac{(n+1)\tilde{\rho}}{mn}}\Big)=\lim_{n\to \infty} n^{1-\tilde{\rho}+\frac{\tilde{\rho}}{m}}= 1$$
Therefore, using Equations~\eqref{TypeCalculation} and~\eqref{GeneralTaylorCoefficient}, combining the computations above we obtain: 
\begin{align*}
    \tau = \frac{1}{e\tilde{\rho}} \limsup_{n \to\infty} (n \sqrt[n]{|c_n|^{\tilde{\rho}}}) &= \frac{1}{e\tilde{\rho}}\lim_{n\to\infty}n\cdot(2\pi)^{\tilde{\rho}}\cdot a^{-\tilde{\rho}/m}\cdot m^{-\tilde{\rho}/m}\cdot n^{-1}\cdot e\\ 
    &= \frac{1}{\tilde{\rho}}  (2\pi)^{\tilde{\rho}}\cdot a^{-\tilde{\rho}/m}\cdot m^{-\tilde{\rho}/m}\\
    &= \frac{m-1}{m}(2\pi)^{m/(m-1)}(am)^{-1/(m-1)}.
\end{align*}
\end{proof}

\end{document}
